\documentclass[reqno]{amsart}
\usepackage[T1]{fontenc}
\usepackage{amsmath,amssymb,amsthm}
\usepackage{xr-hyper}
\usepackage[hidelinks]{hyperref}
\externaldocument[da-]{20261126T030300--the-data__workbook}
\externaldocument[ob-]{20261126T040100--spaces-forms-and-twists__workbook}
\externaldocument[ta-]{20261126T030600--theorem-a-short-forms__workbook}
\externaldocument[dch-]{20261126T030700--the-determinant-of-a-chain__workbook}
\externaldocument[c3-]{20261126T050600--the-exchange-identity__workbook}
\externaldocument[c4-]{20261126T050700--squares-and-the-triangle__workbook}
\externaldocument[tc-]{20261126T040200--theorem-c__workbook}

\newcounter{chap}
% Chapter numbers, generated from 20261125T012000--grh-workbook__guide.org
% by scripts/make-chapters.py.  Do not edit: change the order of the
% index and rerun the script (./preview-site.sh and CI do this).
%
% Usage in a sheet, after \newcounter{chap}:
%   % Chapter numbers, generated from 20261125T012000--grh-workbook__guide.org
% by scripts/make-chapters.py.  Do not edit: change the order of the
% index and rerun the script (./preview-site.sh and CI do this).
%
% Usage in a sheet, after \newcounter{chap}:
%   % Chapter numbers, generated from 20261125T012000--grh-workbook__guide.org
% by scripts/make-chapters.py.  Do not edit: change the order of the
% index and rerun the script (./preview-site.sh and CI do this).
%
% Usage in a sheet, after \newcounter{chap}:
%   \input{chapters}\chapnumber{<sheet ID>}
\makeatletter
\providecommand\chapters@set[2]{\@namedef{chapters@#1}{#2}}
\chapters@set{20261125T001000}{1}% zeros-to-ramp-sums
\chapters@set{20261125T001100}{2}% inert-primes-and-estimate-b
\chapters@set{20261125T005000}{3}% galois-fourier-and-averaging
\chapters@set{20261125T005500}{4}% columns-from-a-sphere
\chapters@set{20261125T000000}{5}% slope-method
\chapters@set{20261126T030100}{6}% a-twist-for-each-prime
\chapters@set{20261125T003000}{7}% the-chain-and-the-square-root
\chapters@set{20261125T010000}{8}% rank-lemma-and-zero-estimate
\chapters@set{20261125T010500}{9}% poincare-series-and-near-orthogonality
\chapters@set{20261124T223000}{10}% elliptic-curves-for-the-proof
\chapters@set{20261125T009500}{11}% fock-model-and-theta-intertwiners
\chapters@set{20261124T223100}{12}% complex-multiplication-and-inert-reduction
\chapters@set{20261125T008000}{13}% one-cm-determinant
\chapters@set{20261124T214500}{14}% supersingular-theta-divisibility
\chapters@set{20261124T214600}{15}% the-isotropic-plane
\chapters@set{20261125T013700}{16}% pairings-on-a-formal-group-of-height-two
\chapters@set{20261125T013500}{17}% hecke-nondegeneracy-on-the-tree
\chapters@set{20261126T030200}{18}% the-bound-is-attained
\chapters@set{20261125T009000}{19}% kronecker-limit-and-chowla-selberg
\chapters@set{20261125T002000}{20}% definite-theta-lifts
\chapters@set{20261125T015500}{21}% siegel-weil-in-the-smallest-case
\chapters@set{20261125T002100}{22}% definite-theta-lifts-and-norms
\chapters@set{20261125T016500}{23}% hecke-operators-and-theta-lifts
\chapters@set{20261125T016700}{24}% divisibility-along-the-chain
\chapters@set{20261125T004000}{25}% denominators-and-sha
\chapters@set{20261125T004100}{26}% tate-twists-and-sha
\chapters@set{20261126T020430}{27}% the-minus-part-by-gauss-sums
\chapters@set{20261125T014100}{28}% congruence-modules
\chapters@set{20261125T014500}{29}% eisenstein-congruences-and-selmer-classes
\chapters@set{20261125T014700}{30}% pseudocharacters-and-the-length-bound
\chapters@set{20261125T011000}{31}% galois-representations-black-box
\chapters@set{20261126T030300}{32}% the-data
\chapters@set{20261126T020210}{33}% bessel-poincare-series-and-their-gram-operators
\chapters@set{20261126T030400}{34}% the-seesaw-identity
\chapters@set{20261126T030500}{35}% divisibility-of-the-composite-lift
\chapters@set{20261126T030600}{36}% theorem-a-short-forms
\chapters@set{20261126T030700}{37}% the-determinant-of-a-chain
\chapters@set{20261126T030800}{38}% the-cost-of-one-lift
\chapters@set{20261126T030900}{39}% denominators-of-projections-to-lifts
\chapters@set{20261126T031000}{40}% twists-with-few-congruences
\chapters@set{20261126T031100}{41}% theorem-c-and-the-main-inequality
\providecommand\chapnumber[1]{%
  \@ifundefined{chapters@#1}%
    {\PackageError{chapters}{Sheet #1 is not listed in the index}%
      {List it in 20261125T012000--grh-workbook__guide.org\MessageBreak
       and run scripts/make-chapters.py.}}%
    {\edef\thechap{\@nameuse{chapters@#1}}}}
\makeatother
\chapnumber{<sheet ID>}
\makeatletter
\providecommand\chapters@set[2]{\@namedef{chapters@#1}{#2}}
\chapters@set{20261125T001000}{1}% zeros-to-ramp-sums
\chapters@set{20261125T001100}{2}% inert-primes-and-estimate-b
\chapters@set{20261125T005000}{3}% galois-fourier-and-averaging
\chapters@set{20261125T005500}{4}% columns-from-a-sphere
\chapters@set{20261125T000000}{5}% slope-method
\chapters@set{20261126T030100}{6}% a-twist-for-each-prime
\chapters@set{20261125T003000}{7}% the-chain-and-the-square-root
\chapters@set{20261125T010000}{8}% rank-lemma-and-zero-estimate
\chapters@set{20261125T010500}{9}% poincare-series-and-near-orthogonality
\chapters@set{20261124T223000}{10}% elliptic-curves-for-the-proof
\chapters@set{20261125T009500}{11}% fock-model-and-theta-intertwiners
\chapters@set{20261124T223100}{12}% complex-multiplication-and-inert-reduction
\chapters@set{20261125T008000}{13}% one-cm-determinant
\chapters@set{20261124T214500}{14}% supersingular-theta-divisibility
\chapters@set{20261124T214600}{15}% the-isotropic-plane
\chapters@set{20261125T013700}{16}% pairings-on-a-formal-group-of-height-two
\chapters@set{20261125T013500}{17}% hecke-nondegeneracy-on-the-tree
\chapters@set{20261126T030200}{18}% the-bound-is-attained
\chapters@set{20261125T009000}{19}% kronecker-limit-and-chowla-selberg
\chapters@set{20261125T002000}{20}% definite-theta-lifts
\chapters@set{20261125T015500}{21}% siegel-weil-in-the-smallest-case
\chapters@set{20261125T002100}{22}% definite-theta-lifts-and-norms
\chapters@set{20261125T016500}{23}% hecke-operators-and-theta-lifts
\chapters@set{20261125T016700}{24}% divisibility-along-the-chain
\chapters@set{20261125T004000}{25}% denominators-and-sha
\chapters@set{20261125T004100}{26}% tate-twists-and-sha
\chapters@set{20261126T020430}{27}% the-minus-part-by-gauss-sums
\chapters@set{20261125T014100}{28}% congruence-modules
\chapters@set{20261125T014500}{29}% eisenstein-congruences-and-selmer-classes
\chapters@set{20261125T014700}{30}% pseudocharacters-and-the-length-bound
\chapters@set{20261125T011000}{31}% galois-representations-black-box
\chapters@set{20261126T030300}{32}% the-data
\chapters@set{20261126T020210}{33}% bessel-poincare-series-and-their-gram-operators
\chapters@set{20261126T030400}{34}% the-seesaw-identity
\chapters@set{20261126T030500}{35}% divisibility-of-the-composite-lift
\chapters@set{20261126T030600}{36}% theorem-a-short-forms
\chapters@set{20261126T030700}{37}% the-determinant-of-a-chain
\chapters@set{20261126T030800}{38}% the-cost-of-one-lift
\chapters@set{20261126T030900}{39}% denominators-of-projections-to-lifts
\chapters@set{20261126T031000}{40}% twists-with-few-congruences
\chapters@set{20261126T031100}{41}% theorem-c-and-the-main-inequality
\providecommand\chapnumber[1]{%
  \@ifundefined{chapters@#1}%
    {\PackageError{chapters}{Sheet #1 is not listed in the index}%
      {List it in 20261125T012000--grh-workbook__guide.org\MessageBreak
       and run scripts/make-chapters.py.}}%
    {\edef\thechap{\@nameuse{chapters@#1}}}}
\makeatother
\chapnumber{<sheet ID>}
\makeatletter
\providecommand\chapters@set[2]{\@namedef{chapters@#1}{#2}}
\chapters@set{20261125T001000}{1}% zeros-to-ramp-sums
\chapters@set{20261125T001100}{2}% inert-primes-and-estimate-b
\chapters@set{20261125T005000}{3}% galois-fourier-and-averaging
\chapters@set{20261125T005500}{4}% columns-from-a-sphere
\chapters@set{20261125T000000}{5}% slope-method
\chapters@set{20261126T030100}{6}% a-twist-for-each-prime
\chapters@set{20261125T003000}{7}% the-chain-and-the-square-root
\chapters@set{20261125T010000}{8}% rank-lemma-and-zero-estimate
\chapters@set{20261125T010500}{9}% poincare-series-and-near-orthogonality
\chapters@set{20261124T223000}{10}% elliptic-curves-for-the-proof
\chapters@set{20261125T009500}{11}% fock-model-and-theta-intertwiners
\chapters@set{20261124T223100}{12}% complex-multiplication-and-inert-reduction
\chapters@set{20261125T008000}{13}% one-cm-determinant
\chapters@set{20261124T214500}{14}% supersingular-theta-divisibility
\chapters@set{20261124T214600}{15}% the-isotropic-plane
\chapters@set{20261125T013700}{16}% pairings-on-a-formal-group-of-height-two
\chapters@set{20261125T013500}{17}% hecke-nondegeneracy-on-the-tree
\chapters@set{20261126T030200}{18}% the-bound-is-attained
\chapters@set{20261125T009000}{19}% kronecker-limit-and-chowla-selberg
\chapters@set{20261125T002000}{20}% definite-theta-lifts
\chapters@set{20261125T015500}{21}% siegel-weil-in-the-smallest-case
\chapters@set{20261125T002100}{22}% definite-theta-lifts-and-norms
\chapters@set{20261125T016500}{23}% hecke-operators-and-theta-lifts
\chapters@set{20261125T016700}{24}% divisibility-along-the-chain
\chapters@set{20261125T004000}{25}% denominators-and-sha
\chapters@set{20261125T004100}{26}% tate-twists-and-sha
\chapters@set{20261126T020430}{27}% the-minus-part-by-gauss-sums
\chapters@set{20261125T014100}{28}% congruence-modules
\chapters@set{20261125T014500}{29}% eisenstein-congruences-and-selmer-classes
\chapters@set{20261125T014700}{30}% pseudocharacters-and-the-length-bound
\chapters@set{20261125T011000}{31}% galois-representations-black-box
\chapters@set{20261126T030300}{32}% the-data
\chapters@set{20261126T020210}{33}% bessel-poincare-series-and-their-gram-operators
\chapters@set{20261126T030400}{34}% the-seesaw-identity
\chapters@set{20261126T030500}{35}% divisibility-of-the-composite-lift
\chapters@set{20261126T030600}{36}% theorem-a-short-forms
\chapters@set{20261126T030700}{37}% the-determinant-of-a-chain
\chapters@set{20261126T030800}{38}% the-cost-of-one-lift
\chapters@set{20261126T030900}{39}% denominators-of-projections-to-lifts
\chapters@set{20261126T031000}{40}% twists-with-few-congruences
\chapters@set{20261126T031100}{41}% theorem-c-and-the-main-inequality
\providecommand\chapnumber[1]{%
  \@ifundefined{chapters@#1}%
    {\PackageError{chapters}{Sheet #1 is not listed in the index}%
      {List it in 20261125T012000--grh-workbook__guide.org\MessageBreak
       and run scripts/make-chapters.py.}}%
    {\edef\thechap{\@nameuse{chapters@#1}}}}
\makeatother

\ifcsname chapters@20261126T060200\endcsname
  \chapnumber{20261126T060200}
\else
  \renewcommand{\thechap}{?}
\fi
\renewcommand{\thesection}{\thechap.\arabic{section}}

\theoremstyle{definition}
\newtheorem{problem}{Problem}[section]
\newtheorem{definition}[problem]{Definition}
\theoremstyle{plain}
\newtheorem{theorem}{Theorem}[section]

\newcommand{\Z}{\mathbb Z}
\newcommand{\R}{\mathbb R}
\newcommand{\C}{\mathbb C}
\newcommand{\F}{\mathbb F}
\newcommand{\Oc}{\mathcal O}
\newcommand{\Nm}{\operatorname{Nm}}
\newcommand{\GL}{\mathrm{GL}}
\newcommand{\Tw}{\mathrm{Tw}}
\newcommand{\cM}{\mathcal M}
\newcommand{\Gb}{\bar G}
\newcommand{\Yt}{\widetilde Y}
\newcommand{\Rm}{\mathfrak R}
\newcommand{\dd}{\mathfrak d}
\newcommand{\ns}{n^\sharp}
\newcommand{\nb}{n^\flat}
\newcommand{\tl}[1]{\texttt{#1}}
\setlength{\emergencystretch}{3em}

\title{Turns and the auxiliary forms}
\author{}
\date{November 26, 2026}

\begin{document}
\maketitle
\refstepcounter{chap}\label{chap}

\paragraph{Prerequisites.}
From Chapter~\ref{da-chap}: the data (Definition~\ref{da-def:data}) and Problem~\ref{da-pr:signs}.  From Chapter~\ref{ob-chap}: the chosen places \(\cM\), the group \(\Gb\), its action \(\rho\) on \(H(Y_0)\) and the idempotent \(e_0\), the target \(T = H_g(Y_0)\), the weights \(\kappa_j\) and the gap \(\Delta\), \(m\)-fold twists \(\vartheta\) with their moduli \(Q_\vartheta\) (we write \(\Tw\) for the set of \(m\)-fold twists), the characters \(\mu_j^0\), \((\omega_w)\) and \(\upsilon\), the lower chain \(\Theta_\vartheta\), partial Hecke values and critical pairs (Definition~\ref{ob-def:pair}), and the target lattice (Definition~\ref{ob-def:target-lattice}).  From Chapter~\ref{ta-chap}, Section~\ref{ta-sec:theorem-a}: integrality of forms at a ring homomorphism \(\eta\colon\C\to\C_p\), the squared norms \(\lvert\sigma F(g)\rvert^2\) and \(\lVert F\rVert_\sigma^2\), and input~3 (the good part).  From Chapter~\ref{c3-chap}: matched functions (Definition~\ref{c3-def:matched}), Theorem~\ref{c3-thm:exchange}, and the nonvanishing of \(L^S(z,\nu)\) for real \(z \ge 1\), finite \(S\) and a character \(\nu\) with a tame part and a weight not identically zero, stated in Section~\ref{c3-sec:integrals}.  From Chapter~\ref{c4-chap}: the relation \(\doteq\) (Definition~\ref{c4-def:doteq}) and Problem~\ref{c4-pr:square-scalars}.  From Chapter~\ref{dch-chap}: Problem~\ref{dch-pr:abs-ratio}.  Finite Fourier analysis on a finite commutative ring.

\medskip
The proof of Theorem~C (Chapter~\ref{tc-chap}) extends a linearly independent family of at least \((1 - 4/b^2)\dim_{\C} T\) target forms to a basis by auxiliary forms.  Their twisted lifts, after multiplication by one scalar \(\gamma(\vartheta)\) for each twist, are integral away from \(\Nm Q_\vartheta\), and their untwisted lifts are short.  Theorem~\ref{thm:aux} expresses \(\lvert\gamma(\vartheta)\rvert\) as a power of \(\Nm Q_\vartheta\) times \(m^2\) turn ratios, each a product of \(t\) ratios of Hecke values; the determinant formula of Chapter~\ref{dch-chap} has \(m(m-1)/2\) pairwise factors.  Why \(m^2\)?  The \(m^2\) factors correspond to the cells crossed in transferring a period through two grids of Hermitian spaces, each crossing using one exchange square.  Sections~\ref{sec:paths}--\ref{sec:one-turn} work out the paths and check one turn; Section~\ref{sec:transfer} states the transfer; Section~\ref{sec:aux} derives Theorem~\ref{thm:aux} from the statement on short lifts and input~3 of Section~\ref{ta-sec:theorem-a}.

\section{Two grids and their paths}\label{sec:paths}

In this sheet \(m\) is even, as for the data; the small values of \(m\) in Problems~\ref{pr:paths} and~\ref{pr:turn-order} concern only the combinatorics of paths.  A \emph{node} is a pair \((i, j)\) with \(0 \le i, j \le m\), and there are two grids of Hermitian spaces over \(E\) on the nodes, with the ranks and signatures specified below and positive signs on all added blocks; the signs refer to the archimedean components of the adelic frames of Section~\ref{ob-sec:spaces}.  The space \(Y^i_j\) of the first grid has rank \(r + (i + j)t\) and \(\lfloor (m-j)/2\rfloor\) negative signs, and \(Y^0_j = Y_j\).  The space \(\Yt^i_j\) of the second grid has rank \(r + 1 + i(t-1) + j(t+1)\) and \(\lfloor (m-i+1)/2 \rfloor\) negative signs, \(\Yt^0_0 = Y_0 \oplus\) a positive line, and on the diagonal, after a choice of basis identification, \(\Yt^j_j = Y^j_j \oplus\) a line of sign \((-1)^j\).  Call the step \((i, j) \to (i, j+1)\) \emph{up} and the step \((i+1, j) \to (i, j)\) \emph{left}.  The matrix of an edge is directed from the end with the smaller coordinate to the end with the larger one, with columns at the former and rows at the latter; its entries are in \(\Rm_\vartheta := \Oc_E/\mathfrak m_\vartheta\) (in this sheet), where
\[
  \mathfrak m_\vartheta := (\pi_M)^b\prod_{w \in Q_\vartheta} w .
\]
Each edge carries a theta lift or a transposed theta lift with the point mass at its matrix as finite datum.  At an end of rank \(d\), the character of an up edge between \(j\) and \(j+1\) is \(\mu^0_{j+1}\vartheta^{j+1}\upsilon^{\lfloor d/2\rfloor}\), and that of a left edge between \(i+1\) and \(i\) is \(\mu^0_{i+1}\upsilon^{\lceil d/2\rceil}\), without twist.  \emph{Status}: read from the proof.

A \emph{path} (in this sheet) goes from \((m, 0)\) to \((0, m)\) by \(m\) steps left (\(L\)) and \(m\) steps up (\(U\)); we write it as a word in \(L, U\).  Its nodes with \(i > j\) carry the spaces \(\Yt^i_j\), those with \(i < j\) the spaces \(Y^i_j\), and its node \((l, l)\) on the diagonal both \(\Yt^l_l\) and \(Y^l_l\): the part of the path before it is its half in the second grid, the part after it its half in the first.  For periods the two halves are read from their outer ends \(\Yt^m_0\) and \(Y_m\) towards the diagonal.  In the first grid the maps are then transposed theta lifts along up edges and theta lifts along left edges, acting on conjugate forms; in the second grid they are theta lifts along up edges and transposed theta lifts along left edges.  Throughout, forms at a node have the archimedean type fixed by the proof: conjugate forms in the first grid and forms in the second.  The \emph{first} path is \(L^mU^m\), through \((0,0)\); the \emph{last} is \(U^mL^m\), through \((m,m)\).  The \emph{cells} of a path are the unit squares \((i, j)\) (named by their lower left corner, \(0 \le i, j < m\)) lying below it.  If a path contains \(LU\) with the node \((i, j)\) between the two steps, replacing \(LU\) by \(UL\) is the \emph{turn} at \((i, j)\); it replaces the node \((i, j)\) by \((i+1, j+1)\).

\begin{problem}\label{pr:nodes}
Formally substituting \(m = 4\) in the signature formulas gives \(2, 1, 1, 0, 0\) negative signs for \(Y^0_0, \dots, Y^0_4\), as in Problem~\ref{da-pr:signs}, and \(0, 1, 1, 2, 2\) for \(\Yt^4_0, \Yt^3_0, \dots, \Yt^0_0\).  In general all spaces on the last path, \(\Yt^m_j\) and \(Y^i_m\), are positive definite, and for \(0 \le j < m\), \(\Yt^{j+1}_j\) and \(Y^{j+1}_j\) have the same rank and the same number of negative signs.
\end{problem}

\begin{problem}\label{pr:paths}
There are \(\binom{2m}{m}\) paths (\(2, 6, 20\) for \(m = 1, 2, 3\)), and each meets the diagonal \(i = j\) in exactly one node \((l, l)\).  The cells of a path form a Young diagram: the column \(i\) contains the cells \((i, 0), \dots, (i, h_i - 1)\), with \(h_0 \ge h_1 \ge \dots \ge h_{m-1}\).
\end{problem}

\begin{problem}\label{pr:cells}
A turn at \((i, j)\) adds exactly the cell \((i, j)\) to the cells of the path.  It is \emph{diagonal} if \(i = j\); then the node \((j, j)\) on the diagonal is replaced by \((j+1, j+1)\).
\end{problem}

\begin{problem}\label{pr:all-cells}
Every path other than the last has a turn.  Any sequence of turns from the first path to the last has exactly \(m^2\) turns, one at each \((i, j)\) with \(0 \le i, j < m\), and \(m\) of them are diagonal.
\end{problem}

\begin{problem}\label{pr:turn-order}
Replacing at each stage the first \(LU\) of the word by \(UL\), the turns from the first path to the last are at \((0,0)\) for \(m = 1\); at \((0,0), (1,0), (0,1), (1,1)\) for \(m = 2\); and at \((0,0), (1,0), (2,0), (0,1), (1,1), (2,1), (0,2), (1,2), (2,2)\) for \(m = 3\).  The numbers of sequences of turns from the first path to the last are \(1, 2, 42\) for \(m = 1, 2, 3\).
\end{problem}

So the paths count \(m^2\) cells, against the \(m(m-1)/2\) squares of the triangle of Chapter~\ref{c4-chap} (Problem~\ref{c4-pr:pairs}): \(1, 4, 9\) against \(0, 1, 3\) for \(m = 1, 2, 3\).

\section{One turn}\label{sec:one-turn}

For a turn at \((i, j)\), call \((i,j)\) the \emph{lost} node, \((i+1,j+1)\) the \emph{gained} node, and \((i, j+1)\) and \((i+1, j)\) the \emph{kept} nodes.  Off the diagonal the square uses the first grid if \(i < j\) and the second if \(i > j\), with the copy of a kept node on the diagonal from that grid; a diagonal turn uses the square at \((j, j)\) in the first grid, \(Y^{j+1}_j\) having the rank and signature of \(\Yt^{j+1}_j\) (Problem~\ref{pr:nodes}).

\begin{definition}\label{def:turn-ratio}
For nonnegative integers \(i, j\), set \(\nu := \mu^0_{i+1}\mu^0_{j+1}\).  For \(\vartheta \in \Tw\) define the \emph{turn ratio} by
\[
  \rho'_{ij}(\vartheta) := \prod_{z=1}^{t} \frac{\lvert L^{\cM \cup Q_\vartheta}\bigl(z, (\nu\vartheta^{j+1})^{-1}\bigr)\rvert}{\lvert L^{\cM}(z, \nu^{-1})\rvert} ,
\]
a quotient by \(0\) being \(0\).
\end{definition}

Here \(\vartheta^{j+1}\) uses the second index.  We use the ratios for \(0 \le i, j < m\).

\begin{problem}\label{pr:turn-pair}
Let \(0 \le i, j < m\).  The character \(\mu^0_{i+1}\mu^0_{j+1}\) has weight \(-(\kappa_{i+1} + \kappa_{j+1})\), with \(t \le \kappa_{i+1}(\tau) + \kappa_{j+1}(\tau) \le 4b\) for every \(\tau\), and tame part \((\omega_w^{i+j+2})\), nontrivial exactly at the \(w \in \cM\).  So each \(P_z := (\mu^0_{i+1}\mu^0_{j+1}, z)\), \(1 \le z \le t\), is a ramified critical pair, \(\rho'_{ij}(\vartheta) = \prod_{z=1}^t R_{P_z}(\vartheta^{j+1}, 1)\), and \(\rho'_{ij}(1) = 1\).
\end{problem}

\begin{problem}\label{pr:turn-ranks}
For \(0 \le i, j < m\), consider the square at \((i, j)\) in either grid, with its lost, gained and kept nodes as in this section.  The gained node has rank \(2t\) more than the lost one; the ranks of the two kept nodes add up to those of the lost and the gained node, and so do their numbers of negative signs; if \(n_0\) is the sum of the ranks of the kept nodes and \(q_-\) the sum of their negative signs, then \(2q_- + 2 \le n_0 \le 2(\kappa_{i+1}(\tau) + \kappa_{j+1}(\tau)) + 2q_-\) for every \(\tau\).  These are the rank and signature conditions of Theorem~\ref{c3-thm:exchange} for \(V^0\) the sum of the kept nodes, \(V^\sharp\) the gained and \(V^\flat\) the lost node, \(\ns - \nb = 2t\), \(h = t\) and \(\mu = \kappa_{i+1} + \kappa_{j+1}\).
\end{problem}

Let \(\nu = \mu^0_{i+1}\mu^0_{j+1}\).  The characters of the two edges at the lost node, taken at any two end ranks with sum \(n_0\), multiply to \(\nu\vartheta^{j+1}\upsilon^{n_0/2}\).  Let \(\chi\) be the inverse of the product of the two characters modulo \(\mathfrak m_\vartheta\) given by their tame parts.  The exchange character \((\nu\vartheta^{j+1})^{-1}\) has weight \(\kappa_{i+1} + \kappa_{j+1}\) and induces \(\chi\).

\emph{The square law of a turn} (read from the proof).  Fix \(0 \le i, j < m\) and one of the two grids.  Take the \emph{source} kept node to be \((i, j+1)\) in the first grid and \((i+1, j)\) in the second, and the other kept node as \emph{target}.  There are \(C_\square \in \C\) (possibly \(0\)), a bijection \(\sigma\) from the coordinates of the lost and gained nodes to those of the source and target, and \(c\colon\Tw\to\C^\times\) with \(c \doteq \rho'_{ij}\), with the following property.  For every \(\vartheta\), every form \(f\) at the source and every linear functional \(f^\vee\) on the forms at the target, let \(L'(\mathbf y)\) be \(f^\vee\) of the transposed lift and then the lift of \(f\) through the lost node, using the two row blocks of \(\mathbf y\), and \(L(\mathbf x)\) be \(f^\vee\) of the lift and then the transposed lift of \(f\) through the gained node, using the two column blocks of \(\mathbf x\); the matrices have entries in \(\Rm_\vartheta\), \(\mathbf y\) with rows at the kept nodes and columns at the lost node, \(\mathbf x\) with rows at the gained node and columns at the kept nodes.  Then \(L'(\mathbf y k) = \chi(\det k)L'(\mathbf y)\) and \(L(h\mathbf x) = \chi(\det h)L(\mathbf x)\) for invertible \(k\), \(h\), and
\[
  \sum_{\mathbf y} \Psi(\mathbf y)L'(\mathbf y) = \frac{C_\square}{c(\vartheta)}\sum_{\mathbf x} \Phi(\mathbf x)L(\mathbf x)
\]
for every pair \((\Phi, \Psi)\) matched for \((\sigma, \chi)\).  The square law uses Theorem~\ref{c3-thm:exchange}, with the conditions of Problem~\ref{pr:turn-ranks}, the critical pair of Problem~\ref{pr:turn-pair}, and the four edges of the square.

\begin{problem}\label{pr:completion}
Let \(A\) be a finite commutative ring, \(\chi\) a character of \(A^\times\), \(n = \nb + \ns\), and \(\Psi\) a complex-valued function on \(n \times \nb\) matrices over \(A\) that vanishes off the first \(\nb\) columns \(k_{\mathcal B}\) of the \(k \in \GL_n(A)\) and satisfies \(\chi(\det h)\Psi(\mathbf y h) = \Psi(\mathbf y)\) for \(h \in \GL_{\nb}(A)\).  Then \(\chi(\det k)\Psi(k_{\mathcal B})\) depends only on the last \(\ns\) rows of \(k^{-1}\), and the function \(\Phi\) equal to it there and to \(0\) elsewhere is matched with \(\Psi\) for \((\sigma, \chi)\) (Definition~\ref{c3-def:matched}), where \(\sigma\) identifies the blocks of sizes \(\nb\), \(\ns\) with the first and the last coordinates of \(A^n\).
\end{problem}

Let \(\mathbf e\) be an additive character of a finite commutative ring \(\Rm\) such that \(x \mapsto \mathbf e(ax)\) is nontrivial for every \(a \ne 0\); such characters exist on \(\Rm_\vartheta\), the product of the rings \(\Oc_E/w^b\), \(w \in \cM\), and \(\Oc_E/w\), \(w \in Q_\vartheta\).  For a function \(F\) on columns \(v \in \Rm^u\) put \(\hat F(v') := \sum_{v \in \Rm^u}\mathbf e(v'v)F(v)\), for rows \(v' \in \Rm^{1\times u}\).

\emph{The line of a diagonal turn} (read from the proof).  At a diagonal turn at \((j, j)\), the old column \(v\) is on the added line of \(\Yt^j_j\) and the new row \(v'\) on the added line of \(\Yt^{j+1}_{j+1}\); put \(u := r + (2j+1)t\), the rank of \(Y^{j+1}_j\).  Let \(\mathfrak a_\vartheta\) be the product of a chosen finite idele generating \(\mathfrak m_\vartheta\) and a fixed finite idele of the Weil data, and \(\mathfrak u_\vartheta := \bar{\mathfrak a}_\vartheta/\mathfrak a_\vartheta\).  The proof writes the theta kernel of the old line, with column \(v\) and archimedean polynomial \(\mathsf P\) (and the finite translation \((1, \mathfrak u_\vartheta^j)\) on the line), as \(c_{\mathrm{line}}(\vartheta)\sum_{v' \in \Rm_\vartheta^{1\times u}}\mathbf e_\vartheta(v'v)\) times the kernel of the new line with row \(v'\) and polynomial \(\mathsf T\mathsf P\) (and \((1, \mathfrak u_\vartheta^{j+1})\) on the line), the sum being taken over these products; here \(\mathsf T\) is a map of archimedean polynomials common to all twists, \(\mathbf e_\vartheta\) an additive character of \(\Rm_\vartheta\), and \(\lvert c_{\mathrm{line}}(\vartheta)\rvert^2(\#\Rm_\vartheta)^u = 1\).  The denominator in the transfer of a diagonal turn is \(c/c_{\mathrm{line}}\), with \(c\) from the square law in the first grid at \((j, j)\).

\begin{problem}\label{pr:line}
Over \(\F_3\) with \(\mathbf e(x) = e^{2\pi \mathrm i x/3}\) and \(u = 1\), the transform of the indicator \(F\) of \(\{1\}\) is \(\hat F(v') = \mathbf e(v')\), with \(\sum\lvert\hat F\rvert^2 = 3\), and transforming back, \(\sum_{v'}\mathbf e(v'v)\hat F(v') = 3F(-v)\), for \(v, v' \in \F_3\).  For a finite commutative ring \(\Rm\) with an additive character \(\mathbf e\) as above, every integer \(u \ge 0\) and every \(F\colon \Rm^u \to \C\), \(\sum_{v' \in \Rm^{1\times u}}\lvert\hat F(v')\rvert^2 = (\#\Rm)^u\sum_{v \in \Rm^u}\lvert F(v)\rvert^2\).
\end{problem}

\begin{problem}\label{pr:line-scalar}
\(\#\Rm_\vartheta = \bigl(\prod_{w \in \cM} q_w\bigr)^b\,\Nm Q_\vartheta\) for every \(\vartheta\).  Fix \(0 \le j < m\), with \(c\) and \(c_{\mathrm{line}}\) as in the line of a diagonal turn.  Then \(c_{\mathrm{line}}^{-1} \doteq 1\) and \(c/c_{\mathrm{line}} \doteq \rho'_{jj}\).
\end{problem}

\begin{problem}\label{pr:roots}
Let \(u, \mathcal C\) be nonnegative integers, \(\Rm\) and \(\mathbf e\) as above, and \(\Psi'(v)\) be a complex-valued function of \(v \in \Rm^u\).  If each value is a sum of at most \(\mathcal C\) roots of unity, then for each row \(v'\) and each root of unity \(\zeta\), \(\zeta\sum_{v \in \Rm^u}\mathbf e(v'v)\Psi'(v)\) is a sum of at most \((\#\Rm)^u\mathcal C\) roots of unity.  For a diagonal turn, \(u = r + (2j+1)t \le 2n_m\).
\end{problem}

\section{The transfer from the first path to the last}\label{sec:transfer}

Fix \(\vartheta \in \Tw\).  A \emph{tuple} on a path assigns to each of its edges a matrix over \(\Rm_\vartheta\) of the size of the edge, directed as in Section~\ref{sec:paths}.  A family \(\mathbf g\) of invertible matrices over \(\Rm_\vartheta\), one at each node, is \emph{admissible} if it lies in the level of the node at the two ends \(Y^0_m = Y_m\) and \(\Yt^m_0\) (trivial modulo \(\mathfrak m_\vartheta\) at \(Y_m\), the kernel of reduction modulo \((\pi_M)^b\) at \(\Yt^m_0\)), and if \(\mathbf g\) at \(\Yt^a_a\) is \(\mathbf g\) at \(Y^a_a\) extended by \(1\) on the line.  These families act on the right on tuples by \((S\cdot\mathbf g)_e := \mathbf g_{e^+}^{-1}S_e\mathbf g_{e^-}\), where \(e^-\) and \(e^+\) are the ends of \(e\) with the smaller and the larger coordinate.  For \(\mathcal C \in \mathbb N\), a \emph{path datum} of bound \(\mathcal C\) (in this sheet) is a function \(\dd\) of the tuples with: \(\dd(S\cdot\mathbf g)\prod_e \chi_e(\det \mathbf g_{e^+}/\det \mathbf g_{e^-}) = \dd(S)\) for admissible \(\mathbf g\), \(\chi_e\) the character of the edge modulo \(\mathfrak m_\vartheta\); \(\dd(S) = 0\) unless the matrices on the up edges have left inverses; and every value is a sum of at most \(\mathcal C\) roots of unity.

In this sheet \(\Gamma\) denotes a path.  For a family \(\dd = (\dd_\vartheta)\) of complex-valued functions of the tuples on \(\Gamma\), with diagonal node \((l, l)\), and a real point \(g\) of \(U(\Yt^l_l)\), define the \emph{period} to be the following function of triples \((\vartheta, \varphi', \varphi)\), where \(\varphi'\) is a conjugate form on \(Y_m\) and \(\varphi\) a form on \(\Yt^m_0\), of arbitrary finite levels.  For each tuple \(S\), the maps along the two halves give \(F_S\) on \(Y^l_l\) and \(\widetilde F_S\) on \(\Yt^l_l\); the period is the sum over \(S\) of \(\dd_\vartheta(S)\) times the integral over the adelic quotient of \(U(Y^l_l)\) of \(F_S(x)\,\widetilde F_S\bigl(\mathrm{diag}(x, \mathfrak u_\vartheta^l)(g, 1)\bigr)\), the second entry acting on the added line at the finite places.  For a function \(\mathcal X\) of triples (in this sheet) we write \(\dd \propto_\Gamma \mathcal X\) if for every \(g\) the period is \(\mathfrak c\,\mathcal X\) for some \(\mathfrak c \in \C\) independent of the triple.

For \(k\colon\Tw\to\mathbb N\) and \(c\colon\Tw\to\C^\times\), we say that \(\Gamma'\) \emph{transfers to} \(\Gamma\) with factors \((k, c)\) if for every bound \(\mathcal C\colon\Tw\to\mathbb N\) and all path data \(\dd'_\vartheta\) on \(\Gamma'\) of bound \(\mathcal C(\vartheta)\) there are path data \(\dd_\vartheta\) on \(\Gamma\) of bound \(k(\vartheta)\mathcal C(\vartheta)\) such that for every \(\mathcal X\), \(\dd \propto_\Gamma \mathcal X\) implies \(\dd' \propto_{\Gamma'} \mathcal X/c\).

\emph{One turn} (read from the proof).  For \(0 \le i, j < m\) there is \(c_{ij}\colon \Tw \to \C^\times\) with \(c_{ij} \doteq \rho'_{ij}\) such that every turn at \((i, j)\) from \(\Gamma'\) to \(\Gamma\) transfers \(\Gamma'\) to \(\Gamma\) with factors \(((\#\Rm_\vartheta)^{2n_m}, c_{ij})\).  With the matrices on the preserved edges fixed, the new datum is matched (Problem~\ref{pr:completion}), for the bijection and character of the square, to the Fourier transform of the old datum in the column of the added line, at each new row; off the diagonal the line is absent and the transform is the identity.  Off the diagonal the square law expresses the old sum as \(C_\square/c(\vartheta)\) times the new sum; on the diagonal the comparison of the column and row periods uses \(\mathsf P\) and \(\mathsf T\mathsf P\) and has the factor \(c_{\mathrm{line}}(\vartheta)C_\square/c(\vartheta)\).

\begin{problem}\label{pr:compose}
If \(\Gamma_1\) transfers to \(\Gamma_2\) with factors \((k, c)\) and \(\Gamma_2\) to \(\Gamma_3\) with \((k', c')\), then \(\Gamma_1\) transfers to \(\Gamma_3\) with \((kk', cc')\).  Hence the first path transfers to the last with factors \(\bigl((\#\Rm_\vartheta)^{2n_m m^2}, c\bigr)\), where \(c := \prod_{0 \le i, j < m}c_{ij} \doteq \prod_{0 \le i,j < m}\rho'_{ij}\).
\end{problem}

\emph{The two ends} (read from the proof).  There are path data on the first path, a real point, a constant and a linear map \(\mathcal E\) (in this sheet) from the forms on \(\Yt^m_0\) of principal level \((\pi_M)^b\) onto \(H(Y_0)\), chosen before the twist and the forms, such that the integral of a conjugate form \(\varphi'\) on \(Y_m\) against the lower chain of \(\mathcal E\varphi\) at the twist \(\vartheta\), before the passage to algebraic forms, is the constant times the period of these data.  On the last path every space is definite (Problem~\ref{pr:nodes}), and the proof expresses its period through a finite sum, over the tuples, of class integrals of a pairing of two chains of algebraic forms.  Transferring by Problem~\ref{pr:compose}, bounding the denominators and the size of that finite sum, and reading the values of \(\Theta_\vartheta\) from such periods by pairing with forms supported on one class, the proof obtains the following statement.

\emph{Short lifts} (read from the proof).  There is \(c\colon\Tw\to\C^\times\) with \(c \doteq \prod_{0 \le i, j < m}\rho'_{ij}\) such that \(H(Y_0)\) is spanned by the \(z\) with: (a) for every \(\vartheta \in \Tw\), every prime \(p \nmid \Nm Q_\vartheta\) and every ring homomorphism \(\eta\colon\C\to\C_p\), \((c(\vartheta)/c(1))\,\Theta_\vartheta z\) is integral at \(\eta\); (b) \(\lvert\sigma(\Theta_1z)(g)\rvert^2 \le \exp(x^{202}b^2)\) for every \(\sigma \in \operatorname{Aut}(\C)\) and every \(g \in G_f(Y_m)\).

\section{Auxiliary forms}\label{sec:aux}

\begin{definition}\label{def:aux}
Let \(\gamma\colon\Tw\to\C\).  A form \(y \in T\) is \emph{auxiliary for \(\gamma\)} if
(i) for every prime \(\ell\), every \(\iota \mid \ell\) and every \(\vartheta \in \Tw\) with \(\gcd(\ell, \Nm Q_\vartheta) = 1\), the form \(\gamma(\vartheta)\Theta_\vartheta y\) lies in \(\mathfrak L_{\vartheta,\iota}\), that is, is integral at \(\iota\); and
(ii) \(\lVert y\rVert_\sigma^2 \le \exp(2x_1^2b^2)\) for every \(\sigma \in \operatorname{Aut}(\C)\), with \(\lVert\cdot\rVert_\sigma\) the length of Definition~\ref{ob-def:target-lattice}, computed from \(\Theta_1 y\).
\end{definition}

In the rest of this section \(c\) is the function of the statement on short lifts, and \(\gamma := c/c(1)\).

\begin{problem}\label{pr:gamma}
\(\gamma(\vartheta) \ne 0\) for all \(\vartheta\), \(\gamma(1) = 1\), and there is \(e_\gamma \in \Z\) with \(\lvert\gamma(\vartheta)\rvert = (\Nm Q_\vartheta)^{e_\gamma/2}\prod_{0 \le i, j < m}\rho'_{ij}(\vartheta)\) for all \(\vartheta \in \Tw\).
\end{problem}

\begin{problem}\label{pr:onto}
The map \(z \mapsto \#\Gb\,\rho(e_0)z\) sends \(H(Y_0)\) onto \(T\).  The image of a spanning set of \(H(Y_0)\) spans \(T\).
\end{problem}

\begin{problem}\label{pr:exponent}
\(x^{99} + x^{202}b^2 \le 2x_1^2b^2\).
\end{problem}

\begin{problem}\label{pr:aux-forms}
If \(z\) satisfies (a) and (b) of the statement on short lifts, then \(\#\Gb\,\rho(e_0)z\) is auxiliary for \(\gamma\).
\end{problem}

Problems~\ref{pr:gamma}--\ref{pr:aux-forms}, the statement on short lifts and input~3 of Section~\ref{ta-sec:theorem-a} prove the following.

\begin{theorem}\label{thm:aux}
There are an integer \(e_\gamma\) and a function \(\gamma\colon\Tw\to\C\), with \(\gamma(\vartheta) \ne 0\) for every \(\vartheta\), \(\gamma(1) = 1\) and
\[
  \lvert\gamma(\vartheta)\rvert = (\Nm Q_\vartheta)^{e_\gamma/2}\prod_{0 \le i, j < m}\rho'_{ij}(\vartheta) \qquad (\vartheta \in \Tw),
\]
such that \(T\) is spanned over \(\C\) by the forms that are auxiliary for \(\gamma\).
\end{theorem}

By Problem~\ref{pr:turn-pair}, \(\prod_{i,j<m}\rho'_{ij}\) is a product of \(m^2t\) ratios \(R_P(\vartheta^n, 1)\) of ramified critical pairs with \(1 \le n \le m\); from this Theorem~\ref{dch-thm:small-ratios} bounds, for a family \(\theta\) with small ratios, the height of \(\iota \mapsto \lvert\gamma(\theta_\iota)\rvert^{-1}\), which Chapter~\ref{tc-chap} uses.  We have proved the bookkeeping of the turns: the path count gives one turn ratio for each of the \(m^2\) cells, and each turn uses an exchange square with \(h = t\); the bound on the number of roots of unity in a value of the datum is at most multiplied by \((\#\Rm_\vartheta)^{2n_m}\) at each turn.  The square law, the line reversal, the two ends and the estimates on the last path are inputs from the proof.

\paragraph{In the Lean formalization.}
\begin{sloppypar}
In the Lean proof \cite{GL} (\tl{grh/\allowbreak src/\allowbreak New/}; the namespace \tl{New.\allowbreak Data} is omitted), Theorem~\ref{thm:aux} is \tl{exists\_\allowbreak span\_\allowbreak isAuxiliaryForm\_\allowbreak eq\_\allowbreak top} and Definition~\ref{def:aux} is \tl{IsAuxiliaryForm}, both in \tl{Determinant/\allowbreak Twist/\allowbreak LongForms.lean}, where \tl{integral\_\allowbreak forms\_\allowbreak long} (Theorem~C) uses them with \tl{HasSmallRatios.\allowbreak heightLe\_\allowbreak inv\_\allowbreak rpow\_\allowbreak mul\_\allowbreak prod\_\allowbreak turnRatio}; \(\rho'_{ij}\) is \tl{turnRatio} in \tl{Determinant/\allowbreak AuxiliaryForms/\allowbreak Turns.lean}.  In the list below, bare filenames are relative to \tl{Determinant/\allowbreak AuxiliaryForms/} under \tl{grh/\allowbreak src/\allowbreak New/}, and paths beginning with \tl{Lib/} are relative to \tl{grh/\allowbreak src/}.  Ranks and signatures: \tl{card\_\allowbreak gridIndex}, \tl{negCount\_\allowbreak gridSigns} (\tl{UpperBound/\allowbreak Objects/\allowbreak Parameters.lean}).  Spaces, moduli and edges: \tl{gridSpace}, \tl{gridSpace\_\allowbreak diag}, \tl{gridModulus}, \tl{gridChar}, \tl{edgeChar}, \tl{liftMove}, \tl{transposedMove}, \tl{GridForm}, \tl{New.\allowbreak GridPath.\allowbreak formMap}, \tl{nodeLevel}, \tl{tuplesOn}, \tl{vectorCountBound}; tuples, admissible families, path data, periods and proportionality: \tl{tupleAction}, \tl{edgeCharProd}, \tl{IsAdmissibleFamily}, \tl{IsPathDatum}, \tl{diagonalPeriodForm}, \tl{gridPeriod}, \tl{IsProportional} (\tl{PathData.lean}).  Paths and turns: \tl{New.\allowbreak TwoGridPath}, with \tl{Turn.\allowbreak squares\_\allowbreak eq}, \tl{exists\_\allowbreak turn}, \tl{turn\_\allowbreak induction} (\tl{GridPaths.lean}, which records both halves as paths from \((0,m)\) to the diagonal and names \(\Yt^i_j\) by \((\mathrm{true}, j, i)\)).  In \tl{Turns.lean}: the character product \tl{gridChar\_\allowbreak mul\_\allowbreak gridChar\_\allowbreak eq\_\allowbreak rankTwist}, the square law \tl{exists\_\allowbreak squareLaw\_\allowbreak grid\_\allowbreak hasAbsUpToNormPower} (from \tl{New.\allowbreak Exchange.\allowbreak one\_\allowbreak square\_\allowbreak identity\_\allowbreak of\_\allowbreak isEdgeFamily} in \tl{Determinant/\allowbreak Exchange/\allowbreak OneSquare.lean}), the critical pair \tl{squarePair}, \tl{isRamified\_\allowbreak squarePair}, the finite characters \tl{edgeChar\_\allowbreak columnEdge\_\allowbreak mul\_\allowbreak rowEdge}, the matched transform \tl{Corner.\allowbreak IsTurn.\allowbreak exists\_\allowbreak isPathDatum} (with \tl{New.\allowbreak Exchange.\allowbreak areMatched\_\allowbreak matchOf} in \tl{Determinant/\allowbreak Chain/\allowbreak SquareLaw.lean}), the relation \tl{TransfersPeriods}, the off-diagonal turn \tl{exists\_\allowbreak transfersPeriods\_\allowbreak off\_\allowbreak diagonal}.  In \tl{DiagonalTurns.lean}: \tl{exists\_\allowbreak columnLineKernel\_\allowbreak eq} (evaluated through \tl{subdiagMap}), \tl{columnPeriod\_\allowbreak eq\_\allowbreak mul\_\allowbreak rowPeriod}, \tl{exists\_\allowbreak transfersPeriods\_\allowbreak diagonal}, \tl{exists\_\allowbreak transfersPeriods\_\allowbreak of\_\allowbreak turn}, \tl{exists\_\allowbreak transfersPeriods\_\allowbreak first\_\allowbreak last}.  The two ends: \tl{exists\_\allowbreak isPathDatum\_\allowbreak first} (\tl{FirstPath.lean}), \tl{lastPathPeriod}, \tl{lastPathChain}, \tl{lastPathPairing}, \tl{isProportional\_\allowbreak lastPathPeriod} (\tl{LastPath.lean}), \tl{exists\_\allowbreak integral\_\allowbreak lowerChainUpTo\_\allowbreak eq\_\allowbreak lastPathPeriod} (\tl{SpanningForms.lean}).  The finite translation \(\mathfrak u_\vartheta\) is \tl{New.\allowbreak Modulus.\allowbreak flipScalar} (\tl{Lib/\allowbreak Weil/\allowbreak DualPair/\allowbreak Flip.lean}).  The estimates and the reading of values: \tl{exists\_\allowbreak mul\_\allowbreak lastPathPeriod\_\allowbreak mem\_\allowbreak embValRing}, \tl{norm\_\allowbreak lastPathPeriod\_\allowbreak le}, \tl{exists\_\allowbreak dualPairing\_\allowbreak eq\_\allowbreak mul\_\allowbreak lastPathPeriod}, \tl{exists\_\allowbreak lowerChain\_\allowbreak integral\_\allowbreak and\_\allowbreak weylNormSq\_\allowbreak le}, \tl{span\_\allowbreak isShort\_\allowbreak eq\_\allowbreak top\_\allowbreak gridLevelForms}; short lifts: \tl{exists\_\allowbreak span\_\allowbreak short\_\allowbreak lifts\_\allowbreak eq\_\allowbreak top} (all in \tl{SpanningForms.lean}).
\end{sloppypar}

\section*{Hints}
\begin{description}
\item[\ref{pr:nodes}] The formulas for the numbers of negative signs.  For the last statement compare \(r + 1 + (j+1)(t-1) + j(t+1)\) with \(r + (2j+1)t\).
\item[\ref{pr:paths}] A path is determined by the positions of its \(m\) letters \(L\).  For the cells, read the height of the path above the column \(i\).
\item[\ref{pr:cells}] Only the height of one column changes.
\item[\ref{pr:all-cells}] A word other than \(U^mL^m\) contains \(LU\).  Count the cells of the first and of the last path.
\item[\ref{pr:turn-order}] For the count, a sequence of turns adds the cells of the \(m \times m\) square one by one, keeping a Young diagram; use the hook length formula.
\item[\ref{pr:turn-pair}] Section~\ref{ob-sec:characters} gives the weights and tame parts of the \(\mu_j^0\).  For the bounds, \(\kappa_j \ge \frac32\Delta\) for \(j \ge 1\), and \(\kappa(\tau) \le b\); \((2m + \frac12)\Delta \le b\) follows from \(\Delta = 2r^8\), \(r \le 3x^5\), \(m \le 2\sqrt x\) and \(x^{10^6} \le x_1 \le b\) (Definition~\ref{da-def:data}); \(\rho'_{ij}(1) = 1\) uses the nonvanishing recalled in Section~\ref{c3-sec:integrals}; for nontriviality, \(0 < i + j + 2 \le 2m < q_2\).
\item[\ref{pr:turn-ranks}] The formulas of Section~\ref{sec:paths}.  For the last inequality, all ranks are at most \(2n_m\), and \(n_m \le \Delta/2\).
\item[\ref{pr:completion}] If \(k\), \(k'\) have the same last rows of the inverse, describe all \(k'\) with the same last rows of the inverse as \(k\).
\item[\ref{pr:line}] Expand \(\lvert\hat F\rvert^2\) and use orthogonality of the characters \(v \mapsto \mathbf e(av)\).
\item[\ref{pr:line-scalar}] Chinese remainder theorem; then Problem~\ref{c4-pr:square-scalars}.
\item[\ref{pr:compose}] Apply the second transfer to the data produced by the first, with the bound \(k\mathcal C\).  Then Problem~\ref{pr:all-cells} and Problem~\ref{c4-pr:square-scalars}.
\item[\ref{pr:gamma}] Problems~\ref{dch-pr:abs-ratio} and~\ref{pr:turn-pair}, with \(\Nm Q_\vartheta\) in place of \(N\).
\item[\ref{pr:onto}] \(T = e_0H(Y_0)\) and \(e_0\) is an idempotent.
\item[\ref{pr:exponent}] \(x^{10^6} \le x_1 \le b\).
\item[\ref{pr:aux-forms}] Input~3 of Section~\ref{ta-sec:theorem-a} with \(s = 1\), \(\xi = \vartheta\) and \(\lambda = \gamma(\vartheta)\), then Problem~\ref{pr:exponent}.  An \(\iota \mid \ell\) induces a ring homomorphism \(\C \to \C_\ell\).  By the remark after Definition~\ref{ob-def:target-lattice}, \(\lVert y\rVert_\sigma^2\) is the weighted sum \(\lVert\Theta_1 y\rVert_\sigma^2\) of Section~\ref{ta-sec:theorem-a}.
\end{description}

\begin{thebibliography}{GL}
\bibitem[GL]{GL} L.~Sernau, \emph{GRH for Dirichlet \(L\)-functions, in Lean~4}, a shorter rewriting of the Lean proof, text by Claude, \url{https://github.com/lsernau/grh-lean}, commit \texttt{72fead1}, readable edition \texttt{grh/}.
\end{thebibliography}

\end{document}
